\chapter{Architectural state}\label{ch:state}

This chapter lists all state visible to software on one core, and the shared state of the machine.

\section{Per-core state}

\begin{table}[htbp]
\centering\small
\caption{Architectural state of one core}\label{tab:state}
\begin{tabularx}{\linewidth}{@{}l l X@{}}
\toprule
\textbf{State} & \textbf{Size} & \textbf{Description}\\\midrule
\reg{x0}--\reg{x31} & $32 \times 64$ bits & Integer registers. \reg{x0} reads as zero; writes to it are discarded.\\
\reg{pc} & 64 bits & Address of the current instruction in the I-LS. Always a multiple of 4 when fetched.\\
\reg{c0}--\reg{c15} & 16 capabilities & Capability registers (Chapter~\ref{ch:caps}).\\
\texttt{CTAB[0..}$\textit{CTAB}$\texttt{-1]} & 256 capabilities (the platform's $\textit{CTAB}$, $\ge 64$) & Capability save table, reached only by \insn{csave}/\insn{crestore} (slots 0--63) and \insn{csavex}/\insn{crestorex} (every slot).\\
\reg{tg0}--\reg{tg15} & 16 tags & Transfer tags: each is busy while a transfer issued with it is in flight.\\
\texttt{XMAP[0..}$\textit{NXMAP}$\texttt{-1]} & 16 entries (the platform's $\textit{NXMAP}$, 4--16) & Overlay map: (valid, space, far address, I-LS address).\\
Inbox & 2 per sender & one FIFO queue of capabilities per sending node, filled by \insn{cgive} and \insn{csend}.\\
Mode & 1 bit & Supervisor or user (\csr{status} bit 0); SPP, the mode \insn{eret} returns to (\csr{status} bit 1).\\
Stop request & 1 bit & \csr{ipend}: set through core control, taken as a \cause{stop} trap in user mode.\\
Pending device reads & up to 4 & Device reads that have issued and wait for their device (Section~\ref{sec:devxfer}); each keeps its tag busy.\\
D-LS & 64 KiB & Data Local Store.\\
I-LS & 64 KiB & Instruction Local Store.\\
CSRs & see Appendix~\ref{app:csrs} & Control and status registers.\\
Halted & 1 bit & Set by \insn{halt} in supervisor mode or by a trap the supervisor cannot take.\\
\bottomrule
\end{tabularx}
\end{table}

\subsection{Integer registers}

The 32 integer registers are 64 bits wide. The assembler accepts \reg{x0}--\reg{x31} and the ABI
names of Table~\ref{tab:abi}; Appendix~\ref{app:abi} gives the calling convention.

\begin{table}[htbp]
\centering\small
\caption{Integer register ABI names}\label{tab:abi}
\begin{tabular}{@{}l l l @{\hspace{3em}} l l l@{}}
\toprule
\textbf{Reg} & \textbf{ABI} & \textbf{Role} & \textbf{Reg} & \textbf{ABI} & \textbf{Role}\\\midrule
\reg{x0} & \reg{zero} & constant zero & \reg{x10}--\reg{x11} & \reg{a0}--\reg{a1} & arguments, results\\
\reg{x1} & \reg{ra} & return address & \reg{x12}--\reg{x17} & \reg{a2}--\reg{a7} & arguments\\
\reg{x2} & \reg{sp} & stack pointer & \reg{x18}--\reg{x27} & \reg{s2}--\reg{s11} & callee-saved\\
\reg{x3} & \reg{gp} & global pointer & \reg{x28}--\reg{x31} & \reg{t3}--\reg{t6} & temporaries\\
\reg{x4} & \reg{tp} & thread pointer & & & \\
\reg{x5}--\reg{x7} & \reg{t0}--\reg{t2} & temporaries & & & \\
\reg{x8} & \reg{s0} / \reg{fp} & callee-saved, frame pointer & & & \\
\reg{x9} & \reg{s1} & callee-saved & & & \\
\bottomrule
\end{tabular}
\end{table}

\subsection{Capability registers and the save table}

The sixteen capability registers \reg{c0}--\reg{c15} each hold either a capability or nothing
(\term{empty}, the null capability). Every far access names one. The \term{capability save
table} gives the supervisor somewhere to keep capabilities while a program owns the registers. Its size
$\textit{CTAB}$ is a property of the platform, given by the platform table's header
(Section~\ref{sec:stpd}): 256 entries by default, at least 64. \insn{csave} and \insn{crestore} name
slots 0--63 with an immediate, \insn{csavex} and \insn{crestorex} any slot with a register. The table is
not addressable memory, so capabilities can never be read or written as data, and user-mode code cannot
reach it. Capabilities are described in Chapter~\ref{ch:caps}.

\subsection{Transfer tags}

Each transfer names one of sixteen tags \reg{tg0}--\reg{tg15}. A tag is \term{busy} from the issue of a
transfer until its completion. Issuing a transfer on a busy tag waits for the tag to become free.
\insn{twait} waits for a set of tags and \insn{tpoll} reports which are busy. Because transfers are
as-if-synchronous, tags affect only \emph{when} a core waits, never what it computes. The tag of a
pending device read is the handle by which a core waits for the device's event: \insn{tany} returns when
the tag becomes idle.

\subsection{Overlay map and inbox}

The overlay map (Section~\ref{sec:overlays}) holds $\textit{NXMAP}$ translations (16 by default, at
least 4: the platform table's header says) from far code addresses (or, with \csr{xslot} set, far code
units aligned to their size) to I-LS addresses for \insn{xcall} and \insn{xret}. The inbox (Section~\ref{sec:inbox}) holds up to two capabilities from each sender, given to
this core, each tagged with the node that sent it. The inbox belongs to the supervisor: the
instructions that take from it or inspect it by sender number are privileged
(Section~\ref{sec:userinbox}).

\section{Control and status registers}\label{sec:csrs}

CSRs are read with \insn{csrr} and written with \insn{csrw}. Table~\ref{tab:csrs} in
Appendix~\ref{app:csrs} defines each one. In summary:
\begin{itemize}
\item counters: \csr{cycle}, \csr{instret}, \csr{stalls};
\item traps, banked by mode: \csr{cause}, \csr{epc}, \csr{tval}, \csr{tvec}; scratch: \csr{scratch};
\item privilege: \csr{status}, the user trap registers \csr{ucause}, \csr{uepc}, \csr{utval},
  \csr{utvec}, and \csr{sscratch}, \csr{ipend} (Chapter~\ref{ch:traps});
\item debug control: \csr{ccr} (bit 0, PTRAP);
\item identity: \csr{coreid}, \csr{ncores};
\item software cache geometry: \csr{ttag}, \csr{tdata}, \csr{tgeom};
\item code paging: \csr{xslot};
\item inter-core: \csr{takefrom};
\item supervised execution: \csr{corebase}, \csr{corecount}.
\end{itemize}
CSR 17 is reserved: it reads as zero and writes to it are ignored.

\section{Shared state}

All cores share T2, T3, the devices and the \term{granule table}: for each 256~KiB granule of T3 (256
of them) and of T2 (one), a 16-bit generation and a head bit (Section~\ref{sec:domains}). Nothing else
is shared: in particular, no core can read another core's registers, Local Stores, save table or
inbox. The only ways to affect another core are far memory (through capabilities), gifts
(\insn{cgive}, \insn{csend}), revoking memory whose capabilities it holds, and a stop request through
core control.

\section{Reset state}\label{sec:reset}

At reset every core starts with:
\begin{itemize}
\item supervisor mode with \csr{status}.SPP = 1 (\csr{status} reads 3);
\item \reg{pc} = 0, all integer registers zero, all other CSRs zero, all tags free, the overlay map and
  inbox empty, and the save table empty; every granule's generation 0 and head bit clear;
\item its I-LS and D-LS loaded from the boot image (the same image on every core; see the STX2 format,
  Section~\ref{sec:stx});
\item the reset capabilities of Table~\ref{tab:resetcaps}. All other capability registers are empty.
\end{itemize}

\begin{table}[htbp]
\centering\small
\caption{Capabilities present at reset}\label{tab:resetcaps}
\begin{tabularx}{\linewidth}{@{}l l l l l X@{}}
\toprule
\textbf{Reg} & \textbf{Cores} & \textbf{Space} & \textbf{Range} & \textbf{Perms} & \textbf{Purpose}\\\midrule
\reg{c0} & core 0 & T3 & $[0, 64~\text{MiB})$ & \perm{RWX} & memory root: all of system memory, the framebuffer window included\\
\reg{c1} & core 0 & T2 & $[0, 256~\text{KiB})$ & \perm{RWX} & all of the cluster store\\
\reg{c2} & every core & device & $[0, 1)$ & \perm{W}\perm{D} & console output port only\\
\reg{c3} & core 0 & device & $[0, \hex{5000})$ & \perm{RWD} & device root: the whole device space\\
\reg{c4} & core 0 & core control & $[0, 8N)$ & \perm{W}\perm{D} & core control: stop any core (Section~\ref{sec:corectl})\\
\bottomrule
\end{tabularx}
\end{table}

The sizes are those of the prototype platform (Section~\ref{sec:memmap}); software reads them from
the platform table. Since \reg{c0} and \reg{c1} are linear and exist only on core 0, every byte of T2 and
T3 starts out owned by exactly one core. Other cores obtain memory only when core 0 gives it to them. No reset
capability is revocable.
